% =============================================================
% MACROS Y BLOQUES REUTILIZABLES
% =============================================================
\newtcolorbox{infobox}[1][]{%
    enhanced,
    before skip=12pt,
    after skip=12pt,
    colback=SoftGray, % Fondo neutro sutil para homogeneidad
    frame hidden,     % Oculta el marco residual
    borderline west={1.2pt}{0pt}{MutedGray}, % Línea gris muy sutil
    arc=1.5mm,
    boxsep=0pt,
    left=12pt,
    right=12pt,
    top=8pt,
    bottom=8pt,
    fonttitle=\bfseries\color{MutedGray}\scshape,
    attach title to upper,
    after title={\par\smallskip},
    title={Información general},
    #1
}

\newtcolorbox{resultbox}[1][]{%
    enhanced,
    before skip=12pt,
    after skip=12pt,
    colback=SoftGray, % Fondo neutro sutil
    frame hidden,
    borderline west={1.2pt}{0pt}{Primary}, % Línea turquesa de énfasis
    arc=1.5mm,
    boxsep=0pt,
    left=12pt,
    right=12pt,
    top=8pt,
    bottom=8pt,
    fonttitle=\bfseries\color{PrimaryDark}\scshape,
    attach title to upper,
    after title={\par\smallskip},
    title={Resultado clave},
    #1
}

\newtcolorbox{warningbox}[1][]{%
    enhanced,
    before skip=12pt,
    after skip=12pt,
    colback=SoftGray, % Fondo neutro sutil
    frame hidden,
    borderline west={1.2pt}{0pt}{MutedGray}, % Línea gris sutil
    arc=1.5mm,
    boxsep=0pt,
    left=12pt,
    right=12pt,
    top=8pt,
    bottom=8pt,
    fonttitle=\bfseries\color{MutedGray}\scshape,
    attach title to upper,
    after title={\par\smallskip},
    title={Nota de simulación},
    #1
}
\newcommand{\captura}[3][0.88]{%
\begin{figure}[H]
    \centering
    \includegraphics[width=#1\linewidth]{#2}
    \caption{#3}
\end{figure}}
\newcommand{\dato}[2]{\textbf{#1:} #2\\}

% -------------------------------------------------------------
% NUEVAS HERRAMIENTAS DE DISEÑO SUBTILES Y FORMALES
% -------------------------------------------------------------

% 1. Editor de Código estilo formal (Pestaña limpia)
\newtcolorbox{codeeditor}[2][]{%
    enhanced,
    boxrule=0.4pt,
    colframe=RuleGray,
    colback=SoftGray,
    arc=1.5mm,
    top=14pt, % Margen superior optimizado sin botones de colores
    bottom=6pt,
    left=6pt,
    right=6pt,
    fonttitle=\bfseries\small\ttfamily,
    coltitle=MutedGray,
    colbacktitle=SoftGray,
    attach boxed title to top left={yshift=-3mm, xshift=5mm},
    boxed title style={
        enhanced,
        boxrule=0.4pt,
        colframe=RuleGray,
        colback=SoftGray,
        arc=1mm,
        top=2pt,
        bottom=2pt,
        left=5pt,
        right=5pt,
    },
    title={#2},
    underlay={
        % Línea de separación horizontal de la cabecera gris ultra fina
        \draw[RuleGray, line width=0.3pt] ([yshift=-7mm]frame.north west) -- ([yshift=-7mm]frame.north east);
    },
    #1
}

\newtcolorbox{conceptcard}[2][]{%
    enhanced,
    before skip=12pt,
    after skip=12pt,
    colback=SoftGray, % Fondo sutil para detectarlo rápido al deslizar la vista
    frame hidden,     % Oculta completamente el marco residual
    borderline west={1.2pt}{0pt}{Primary}, % Línea de énfasis turquesa muy sutil
    arc=1.5mm,        % Esquinas redondeadas y modernas
    boxsep=0pt,
    left=12pt,  % Margen izquierdo respecto a la línea
    right=12pt, % Margen derecho para contener el texto
    top=8pt,    % Margen superior interno
    bottom=8pt, % Margen inferior interno
    fonttitle=\bfseries\color{PrimaryDark}\scshape,
    attach title to upper,
    after title={\par\smallskip},
    title={#2}, % Título formal y limpio, sin iconos decorativos
    #1
}

% 3. Tarjeta con Código QR Vectorial (4 parámetros: 1 opcional + 3 obligatorios)
\newcommand{\enlaceQR}[4][1.2cm]{%
\begin{tcolorbox}[
    enhanced,
    colback=SoftGray,
    colframe=RuleGray,
    boxrule=0.4pt,
    arc=1mm,
    width=\linewidth,
    left=8pt,
    right=8pt,
    top=6pt,
    bottom=6pt,
]
\begin{minipage}[c]{0.70\linewidth}
    \small\textbf{\color{PrimaryDark}\scshape #3}\\ % Título en versalitas
    \vspace{2pt}
    \scriptsize\color{TextGray}#4\\ % Descripción
    \vspace{2.5pt}
    \tiny\url{#2} % URL del recurso
\end{minipage}\hfill
\begin{minipage}[c]{0.26\linewidth}
    \centering
    \qrcode[height=#1]{#2} % Código QR vectorial sutil
\end{minipage}
\end{tcolorbox}}
